\section{A signed kernel for polynomial differences}
\label{sec:kernel}

The tuple laws constrain differences modulo primes. We now connect their
pair distributions to actual values of $h_l$ at positive integers. The
connection is a uniform Fourier approximation: a kernel supported on those
polynomial values will have the same small-denominator multipliers as the
pair law, up to a power error. We adapt the signed square-kernel construction
of~\cite[Section~6, especially Proposition~6.1]{squarepower}. For a general
polynomial, we realize the pair distribution by weighting regular argument
classes; its explicit probability formula is not needed.

\subsection{Pair multipliers and the ordered estimate}

Fix an integer $l\ge1$, put $F=h_l$, and choose a directed edge
$v_j\to v_{j+1}$ of the designated cycle. For $p\ge P$, let $\pi_p$
be the distribution of $z_{v_{j+1}}-z_{v_j}$ under $\mu_{l,p}^{\circ}$.
Write $\pi_p(y)$ for its point probabilities and define
\[
 \Gamma_p(a'/p)=\sum_{y\in\F_p}\pi_p(y)e(a'y/p).
\]
The distribution is supported on values of $F$ at regular arguments.
Conditioning the character of the second endpoint on the other slots,
then pairing with the opposite character of the first endpoint, gives
\[
 |\Gamma_p(a'/p)|\le p^{-\gamma}\qquad(p\nmid a')
\]
by~\eqref{eq:8} and Cauchy--Schwarz. For a rational $\xi$ with squarefree
denominator supported on primes at least $P$, decompose $\xi$ additively
into its prime-denominator parts and define $\Gamma(\xi)$ as the product
of the corresponding $\Gamma_p$. In particular,
\begin{equation}
 \Gamma(0)=1,\qquad |\Gamma(\xi)|\le q(\xi)^{-\gamma}.
 \label{eq:15}
\end{equation}
All estimates below are uniform in the selected cycle edge.

\begin{proposition}[Ordered polynomial-kernel estimate]
\label{kernel:estimate}
Use the fixed parameters in~\eqref{eq:14}, and let $N$ be a sufficiently
large integer. Suppose $1\le l\le N^\zeta$ and
\[
 D=M_0\prod_{p\in B}p\le N^\beta,
\]
where $B$ is a finite set of primes at least $P$. Let
$J_1,J_2:[N]\to\C$ satisfy $|J_1|,|J_2|\le1$ and be supported in
residue classes $a_1,a_2$ modulo $D$, respectively. Suppose that
$r=a_2-a_1$ is admissible for $l$ modulo $D$, that $x<y$ whenever
$x\in\operatorname{supp}J_1$ and $y\in\operatorname{supp}J_2$, and
that no such difference $y-x$ belongs to $F(\N_+)$.
Put $H=N^{\beta/2}$ and
\[
 \Xi=\{\xi\in\Q/\Z:q(\xi)\le H,\quad q(\xi)\text{ squarefree},
                         \quad(q(\xi),D)=1\}.
\]
Then
\begin{equation}
 D\left|\sum_{\xi\in\Xi}\Gamma(\xi)
          \widehat{J_1}_{[N]}(-\xi)\widehat{J_2}_{[N]}(\xi)\right|
 \ll H^{-\gamma/2}.
 \label{eq:16}
\end{equation}
The implied constant and the threshold for $N$ depend only on the fixed
polynomial and the fixed parameters, not on $l,D,r,J_1,J_2$.
\end{proposition}

Every prime below $P$ divides $D$, so $\Gamma$ is defined at every
frequency in $\Xi$. We prove the proposition by constructing two kernels:
one encodes its Fourier expression, and the other is supported on forbidden
positive polynomial shifts. We first establish the uniform Fourier
comparison; a bilinear identity then gives~\eqref{eq:16}.

\subsection{The model kernel and a kernel on polynomial values}

Fix the data of Proposition~\ref{kernel:estimate}. Define the model kernel
on $\Z$ by
\[
 k_0(n)=\frac DN\1_{\{1\le n\le N\}}\1_{\{n\equiv r\ (D)\}}
                   \sum_{\xi\in\Xi}\Gamma(\xi)e(-\xi n).
\]
We first choose periodic argument weights that will reproduce its
coefficients. For every $p\mid D$, admissibility and the lifting construction
in Section~\ref{sec:local} provide an integer $x_p$ with
$F(x_p)\equiv r\pmod{p^{E_p}}$ and
$s_p=v_p(F'(x_p))$ satisfying $E_p\ge2s_p+1$.
On the argument class $x_p\pmod{p^{E_p-s_p}}$, the polynomial $F$
is constant $r$ modulo $p^{E_p}$ and uniform among its higher lifts.
Combine these classes by the Chinese remainder theorem. Their modulus
$d_r=\prod_{p\mid D}p^{E_p-s_p}$ divides $D$; let $\rho_D$ be $d_r$
times the indicator of the resulting argument class. Thus $\rho_D$ has
periodic mean one and supremum at most $D$.

For $p\nmid D$ (and hence $p\ge P$), choose, for each $y$ with
$\pi_p(y)>0$, one regular
argument $x_y$ satisfying $F(x_y)=y$ in $\F_p$. Define
\[
 \psi_p(x_y)=p\pi_p(y),\qquad \psi_p(x)=0
       \quad\text{at the other arguments}.
\]
Then $\psi_p$ is nonnegative, has uniform mean one, and pushes forward
under $F$ to $\pi_p$. In particular $0\le\psi_p\le p$. For positive
squarefree integers $u$ coprime to $D$, define periodic functions
\[
 B_p(m)=\psi_p(m)-1,\qquad
 B_u(m)=\prod_{p\mid u}B_p(m),\qquad B_1=1.
\]
Each $B_p$ has mean zero, and $|B_u|\le u$. We will use the truncated
sum of these products to select the prescribed local distributions at
small denominators. Indeed, for positive squarefree $R$ coprime to $D$,
\[
 \sum_{u\mid R}B_u(m)=\prod_{p\mid R}(1+B_p(m))
                    =\prod_{p\mid R}\psi_p(m).
\]
The cutoff $u\le H$ preserves this identity when $R\le H$, and we
will use local exponential-sum cancellation for larger denominators.
The truncated sum may have either sign.

All coefficients of $F$ are $O_h(a_l)$. Hence one integer $x_*\ge1$,
depending only on $h$, can be chosen for the entire auxiliary family,
independently of $l$, so that for $x\ge x_*$,
\[
 F(x)\asymp a_lx^k,\qquad F'(x)\asymp a_lx^{k-1},\qquad
 |F''(x)|\ll a_lx^{k-2},
\]
with $F'(x)>0$. The lower-degree terms have total absolute value
$O_h(a_lx^{k-1})$ for $x\ge1$. Enlarging this same $x_*$ therefore also
arranges
\[
 F(x)\ge \frac{a_lx^k}{2}>0\qquad(x\ge x_*),
\]
uniformly in $l$. Put $U=(N/a_l)^{1/k}$. Since $k\zeta=\eta/20$,
the hypothesis $l\le N^\zeta$ gives
\[
 k!a_l\le k!a l^k\le k!a N^{\eta/20}.
\]
The exponent gap $\eta/8-\eta/20=3\eta/40$ is positive. Thus one
sufficiently large threshold for $N$, depending on $h$ and the fixed
parameters but independent of $l$, ensures $k!a\le N^{3\eta/40}$ and hence,
simultaneously for all $l\le N^\zeta$,
\[
 k!a_l\le N^{\eta/8},\qquad
 N^{1-\eta/8}\le U^k\le N.
\]
The last inequality uses that $a_l$ is a positive integer. The lower bound
for $U$ tends to infinity uniformly in this range of $l$, so increasing the
same threshold ensures that there is a unique real $X>x_*$ with $F(X)=N$,
and $X\asymp U$. The lower bound on $F(X)$ gives
\[
 X\le (2N/a_l)^{1/k}=2^{1/k}U<2U.
\]
With $m$ ranging over integers and $u$ over positive integers, define
\begin{equation}
 k_1(n)=\sum_{x_*<m\le X}\1_{\{n=F(m)\}}\frac{F'(m)}N\rho_D(m)
             \sum_{\substack{u\le H\ \mathrm{squarefree}\\(u,D)=1}}B_u(m).
 \label{eq:17}
\end{equation}
This kernel is supported on positive
values in $F(\N_+)$. Derivative weights also occur in Lucier's
polynomial-difference argument~\cite[Section~2]{Lucier2006}.
Here the factor $F'(m)/N$ is chosen so that the
corresponding continuous measure becomes $dt/N$ under $t=F(x)$;
it will therefore match the interval average in the model kernel.

Use positive-sign transforms for the kernels, and write
\[
 S_i(\alpha)=\sum_n k_i(n)e(\alpha n),\qquad
 V_N(y)=\frac1N\sum_{n=1}^N e(yn).
\]
Our main analytic claim is the following.

\begin{lemma}[Uniform kernel comparison]
\label{kernel:comparison}
For the kernels just constructed,
\begin{equation}
 \sup_{\alpha\in\R/\Z}|S_1(\alpha)-S_0(\alpha)|
       \ll H^{-\gamma/2},
 \label{eq:18}
\end{equation}
uniformly in the data of Proposition~\ref{kernel:estimate}.
\end{lemma}

To describe the model transform, expand its congruence indicator:
\begin{equation}
 S_0(\alpha)=\sum_{\xi\in\Xi}\sum_{j\bmod D}
           \Gamma(\xi)e(-jr/D)V_N(\alpha-\xi+j/D).
 \label{eq:19}
\end{equation}
There are at most $DH^2$ terms, all with coefficients of modulus at
most one. Their centers $\xi-j/D$ have denominators at most $DH$ and
are distinct: equality of two centers makes $\xi-\xi'$ have denominator
both coprime to $D$ and dividing $D$, so it is zero.

Every $\alpha_0\in\Q/\Z$ has a unique decomposition
$\alpha_0=\alpha_D+\alpha_1$ in which $q(\alpha_D)$ has only prime
factors dividing $D$ and $(q(\alpha_1),D)=1$. Prime powers are allowed
in both parts. The coefficient of the center $\alpha_0$ in~\eqref{eq:19},
taken to be zero when that center is absent, is
\begin{equation}
 c_0(\alpha_0)=
 \begin{cases}
 e(\alpha_Dr)\Gamma(\alpha_1),&q(\alpha_D)\mid D,
                                      \quad\alpha_1\in\Xi,\\
 0,&\text{otherwise}.
 \end{cases}
 \label{eq:20}
\end{equation}
Indeed, at a center one has $\alpha_1=\xi$ and $\alpha_D=-j/D$.

\subsection{Minor arcs: a uniform Weyl estimate}

Take the major arcs to be the closed balls of radius $3N^{-1+\eta}$
about the rational points with reduced denominator at most $N^\eta$.
They are disjoint for large $N$, since $3\eta<1$. Their complement
is the set of minor arcs. We write $\|t\|_{\R/\Z}$ for distance to
the nearest integer. The periodic weights in $k_1$ will introduce
linear twists, so the required exponential-sum estimate includes them.

\begin{lemma}
\label{kernel:weyl}
If $\alpha$ lies on the minor arcs, then, for every real $y'$ and every
integer interval $I\subseteq(x_*,X]$,
\begin{equation}
 \left|\sum_{m\in I}e\bigl(\alpha F(m)+y'm\bigr)\right|
       \ll U N^{-\nu}.
 \label{eq:21}
\end{equation}
\end{lemma}

\begin{proof}
Put $A_l=k!a_l$, $\theta'=A_l\alpha$, and
$Q_0=\lfloor N^{1-\eta}\rfloor$. Pigeonholing the $Q_0+1$ points
$0,\theta',\ldots,Q_0\theta'$ into $Q_0$ equal intervals gives a reduced
rational $a'/q'$ such that
\[
 1\le q'\le Q_0,\qquad
 \|\theta'-a'/q'\|_{\R/\Z}\le\frac1{q'Q_0}\le\frac1{(q')^2}.
\]
Lifting this approximation through multiplication by $A_l$ places
$\alpha$ within $1/(A_lq'Q_0)$ of a rational with denominator at most
$A_lq'$. If $A_lq'\le N^\eta$, this would place $\alpha$ in a major
arc, since $Q_0^{-1}\le2N^{-1+\eta}$. Therefore
$A_lq'>N^\eta$, and in particular
\[
 N^{\eta/2}\le q'\le N^{1-\eta},
\]
using $A_l\le N^{\eta/8}$.

Choose $M=\lceil2U\rceil$. Since $X<2U$, this integer is larger than all
the interval lengths under consideration. Also $U\ge1$, so the absolute
comparison $2U\le M\le2U+1\le3U$ holds. Repeated differencing gives
\begin{equation}
 \begin{split}
 \left|\sum_{m\in I}e\bigl(\alpha F(m)+y'm\bigr)\right|^{2^{k-1}}
 &\le C_k M^{2^{k-1}-k}\\
 &\quad\times\sum_{|b_1|,\ldots,|b_{k-1}|<M}
 \min\!\left(M,\|\theta'b_1\cdots b_{k-1}\|_{\R/\Z}^{-1}\right).
 \end{split}
 \label{eq:22}
\end{equation}
Here the minimum equals $M$ when the distance is zero. To verify the
estimate, square the sum and group pairs by their difference. After
taking absolute values, square again and use Cauchy--Schwarz on the
increment indices. At each step the inner sum is over an intersection
of interval translates, hence an interval. After $s$ steps the prefactor
is $M^{2^s-s-1}$: its exponent is zero for $s=1$, and the next step
doubles that exponent and adds $s$. After $k-1$ steps the phase is
linear with coefficient $k!a_l\alpha b_1\cdots b_{k-1}$, and a geometric
sum proves~\eqref{eq:22}. The added linear term contributes only a
constant after the first differencing and does not change the estimate.
This derivation uses only $|I|\le M$, so also covers short intervals.

The tuples with a zero increment contribute $O_k(M^{k-1})$ inside
the sum in~\eqref{eq:22}. For the others, group by the absolute product
$m\le M^{k-1}$. For each $\epsilon>0$ their multiplicity is
$O_{\epsilon,k}(m^\epsilon)$. Indeed, at a prime power $p^a$ the
number of exponent choices is a binomial coefficient of fixed degree
in $a$; it is at most $p^{a\epsilon}$ for all sufficiently large $p$.
The finitely many smaller primes contribute a bounded extra factor,
since a fixed polynomial in $a$ divided by $p^{a\epsilon}$ is bounded.

Partition the positive product indices into blocks of
$\lfloor q'/2\rfloor$ consecutive integers. If $0<b<q'/2$, then
\[
 \|\theta'b\|_{\R/\Z}
 \ge\|a'b/q'\|_{\R/\Z}-\frac b{(q')^2}
 \ge\frac1{2q'}.
\]
Thus the points in a block have spacing at least $1/(2q')$ on the
circle. Ordering their distances to zero and summing the reciprocal-distance
bound gives $O(M+q'\log(2q'))$ per block. There are
$O(M^{k-1}/q'+1)$ blocks. Dividing~\eqref{eq:22} by
$M^{2^{k-1}}$, and absorbing the divisor bound and logarithms into
$N^\epsilon$, gives
\[
 C_{\epsilon,k}N^\epsilon
       \left(\frac1{q'}+\frac1M+\frac{q'}{M^k}\right).
\]
Here $M^k\ll N$ and $M^k\gg N^{1-\eta/8}$. The three terms are
bounded respectively by $N^{-\eta/2}$,
$O(N^{-(1-\eta/8)/k})$, and $O(N^{-7\eta/8})$.
For sufficiently small fixed $\epsilon$, the displayed expression
is $O(N^{-\eta/4})$. Taking its $2^{k-1}$-st root proves~\eqref{eq:21},
because $\eta/(4\cdot2^{k-1})=5\nu$.
\end{proof}

We can now bound both kernels on the minor arcs. The supremum and
total variation of $F'/N$ on $[x_*,X]$ are $O(1/U)$, so partial
summation in~\eqref{eq:21} gives $O(N^{-\nu})$ with this weight.
With Fourier coefficients normalized by their period, $\rho_D$ has
Fourier $\ell^1$ norm $d_r\le D$. The corresponding norm of $B_u$
is at most $u^2$, by its period and supremum. Expanding these periodic
weights into linear twists and summing over $u\le H$ yields
\[
 |S_1(\alpha)|\ll DH^3N^{-\nu}.
\]
Every center in~\eqref{eq:19} has denominator at most
$DH\le N^{3\beta/2}\le N^\eta$. On the minor arcs its distance
from $\alpha$ is at least $3N^{-1+\eta}$. The geometric-sum bound
$|V_N(t)|\le\min(1,(2N\|t\|_{\R/\Z})^{-1})$ therefore gives
\[
 |S_0(\alpha)|\ll DH^2N^{-\eta}.
\]
These bounds prove the required comparison away from the major arcs.
Near a small rational, cancellation in the argument variable must be
replaced by an exact calculation of its periodic mean.

\subsection{Major arcs and the local coefficient calculation}

Write $\alpha=\alpha_0+y$, where
$q=q(\alpha_0)\le N^\eta$ and $|y|\le3N^{-1+\eta}$. A model center
other than $\alpha_0$ is at distance at least $1/(qDH)$ from
$\alpha_0$. Since $|y|qDH=o(1)$ uniformly,~\eqref{eq:19} gives
\[
 S_0(\alpha)=c_0(\alpha_0)V_N(y)+O(qD^2H^3/N).
\]

For squarefree $u\le H$ coprime to $D$, let $A_u$ be the periodic
mean of
\[
 b_u(m)=\rho_D(m)B_u(m)e(\alpha_0F(m)).
\]
Its period is at most $qDu$ and its supremum at most $Du$.
Consequently its sum over any real interval differs from $A_u$ times
the interval's length by $O(qD^2u^2)$. For
$w(x)=F'(x)e(yF(x))/N$, differentiation gives
\[
 \|w\|_\infty+\int_{x_*}^{X}|w'(x)|\,dx
 \ll\frac{1+|y|N}{U}.
\]
Partial summation, followed by the substitution $t=F(x)$, gives
\[
 \sum_{x_*<m\le X}\frac{F'(m)}N e(yF(m))b_u(m)
 =\frac{A_u}{N}\int_{F(x_*)}^{N}e(yt)\,dt
    +O\!\left(qD^2u^2\frac{1+|y|N}{U}\right).
\]
The normalized integral differs from $V_N(y)$ by
$O(a_l/N+(1+|y|N)/N)$. Multiplying by $|A_u|\le Du$ leaves an error
covered by the one already displayed, since $a_l/N=U^{-k}$ and
$U\le N$. Summing over $u$ therefore proves
\begin{equation}
 \begin{split}
 S_1(\alpha)&=c(\alpha_0)V_N(y)
       +O\!\left(qD^2H^3\frac{1+|y|N}{U}\right),\\
 c(\alpha_0)&=\sum_{\substack{u\le H\ \mathrm{squarefree}\\(u,D)=1}}A_u.
 \end{split}
 \label{eq:23}
\end{equation}
The analytic approximations now have the same factor $V_N(y)$.
It remains to compare their coefficients. This is where the regular
argument weights $\psi_p$ reproduce the pair law.
The periodic means $A_u$ and their sum $c(\alpha_0)$ are defined by
the same formulas for every rational $\alpha_0$, without a denominator
restriction; we use that domain in the next lemma.

\begin{lemma}[Comparison of rational coefficients]
\label{kernel:coefficients}
For every $\alpha_0\in\Q/\Z$, write $\alpha_0=\alpha_D+\alpha_1$ for the
unique decomposition in which $q(\alpha_D)$ has only prime factors dividing
$D$ and $(q(\alpha_1),D)=1$. Then the coefficients defined by
\eqref{eq:20} and~\eqref{eq:23} satisfy
\[
 |c(\alpha_0)-c_0(\alpha_0)|\le H^{-\gamma/2}.
\]
In fact they agree when $q(\alpha_1)\le H$.
\end{lemma}

\begin{proof}
In a periodic mean we may average at any common multiple of the
periods and separate prime-power factors by the Chinese remainder
theorem. First consider a prime $p\mid D$, and abbreviate
$E=E_p$, $s=s_p$. On the chosen argument class, Section~\ref{sec:local}
gives
\[
 F(x_p+p^{E-s}z)=F(x_p)+p^E(c_1z+pC(z)),\qquad p\nmid c_1,
\]
where $C$ has integer coefficients. The polynomial in parentheses
is a bijection modulo every power of $p$: its difference at two
arguments is their difference times a unit. Thus for a local frequency
of denominator $p^j$ with $j\le E$, the phase is constant at its
value on $r$. This remains true if $j>E-s$. For $j>E$, uniform
arguments in the restricted class modulo $p^j$ give uniform $z$
modulo $p^{j-E+s}$, hence uniform reduction modulo $p^{j-E}$.
The $F$-values are consequently uniform among the lifts of $r$
modulo $p^j$, and a primitive character has mean zero.
The joint factor from primes dividing $D$ is therefore
\[
 \chi_D(\alpha_0)=
 \begin{cases}
 e(\alpha_Dr),&q(\alpha_D)\mid D,\\
 0,&q(\alpha_D)\nmid D.
 \end{cases}
\]

Put $q_1=q(\alpha_1)$. If $u$ has a prime not dividing $q_1$,
its factor $B_p$ has mean zero independently of all the other
factors, so $A_u=0$. Thus only divisors of
$\operatorname{rad}(q_1)=\prod_{p\mid q_1}p$ contribute.
For $p^j\Vert q_1$, let $\alpha_p$ be the primitive $p^j$-denominator
part of $\alpha_1$ and set
\[
 a_{p,j}=\E_{m\bmod p^j}e(\alpha_pF(m)),\qquad
 b_{p,j}=\E_{m\bmod p^j}B_p(m)e(\alpha_pF(m)).
\]
The Chinese remainder theorem now expresses the truncated coefficient as
\begin{equation}
 c(\alpha_0)=\chi_D(\alpha_0)
 \sum_{\substack{u\mid\operatorname{rad}(q_1)\\u\le H}}
 \left(\prod_{\substack{p^j\Vert q_1\\p\mid u}}b_{p,j}\right)
 \left(\prod_{\substack{p^j\Vert q_1\\p\nmid u}}a_{p,j}\right).
 \label{kernel:truncated-coefficient}
\end{equation}
The mean against $1+B_p=\psi_p$ is
\[
 a_{p,j}+b_{p,j}=
 \begin{cases}
 \Gamma_p(\alpha_p),&j=1,\\
 0,&j\ge2.
 \end{cases}
\]
For $j=1$ this is the defining pushforward of $\psi_p$. For
$j\ge2$, condition on a regular argument modulo $p$. Its lifts
are mapped uniformly to the lifts of its value, so their primitive
character mean vanishes. This holds on each argument class separately,
regardless of the size of $\psi_p$ there.

If $q_1\le H$, all divisors of $\operatorname{rad}(q_1)$ occur.
Thus~\eqref{kernel:truncated-coefficient} factors as
\[
 c(\alpha_0)=\chi_D(\alpha_0)
                    \prod_{p^j\Vert q_1}(a_{p,j}+b_{p,j}).
\]
For squarefree $q_1$ this is
$\chi_D(\alpha_0)\Gamma(\alpha_1)$; if any exponent is at least
two, it is zero. These are precisely the cases in~\eqref{eq:20}.
The empty product, when $q_1=1$, also gives the required coefficient.

Suppose instead that $q_1>H$. The model coefficient is then zero.
By~\eqref{eq:2}, $|a_{p,j}|\le p^{-8\delta j}$. For $j=1$ the
preceding identity and~\eqref{eq:15} give
$|b_{p,1}|\le2p^{-\gamma}$; for $j\ge2$ the identity gives
$b_{p,j}=-a_{p,j}$, hence again
$|b_{p,j}|\le2p^{-\gamma j}$. Taking absolute values
in~\eqref{kernel:truncated-coefficient} and dropping the divisor cutoff yields
\begin{align*}
 |c(\alpha_0)|
 &\le\prod_{p^j\Vert q_1}(|a_{p,j}|+|b_{p,j}|)\\
 &\le3^{\omega(q_1)}q_1^{-\gamma}
 \le q_1^{-\gamma/2}<H^{-\gamma/2}.
\end{align*}
Here $\omega(q_1)$ is the number of distinct prime divisors, and the
last inequality before the strict one follows from~\eqref{eq:9},
since every prime of $q_1$ is at least $P$. This bound also covers
$\operatorname{rad}(q_1)\le H<q_1$: in that case the full divisor
product is present and is actually zero, because some exponent is
at least two. If $q(\alpha_D)\nmid D$, all coefficients vanish
already through $\chi_D(\alpha_0)$.
\end{proof}

\begin{proof}[Proof of Lemma~\ref{kernel:comparison}]
On the minor arcs the two bounds already proved give
\[
 |S_1|\ll N^{-\nu+(5/2)\beta},\qquad
 |S_0|\ll N^{-\eta+2\beta}.
\]
On a major arc, the error terms in the two approximations are at most
\[
 O(N^{\eta+4\beta-1}),\qquad
 O\bigl(N^{2\eta+4\beta-(1-\eta/8)/k}\bigr),
\]
respectively. Each of these four bounds is
$O(N^{-\beta\gamma/4})=O(H^{-\gamma/2})$ by~\eqref{eq:14}.
Indeed $\nu=100\beta$, $\eta=1/(100k)$ and
$\beta=\eta/(1000\cdot2^k)$, so the two major-arc exponents are
less than $-0.99$ and $-0.97/k$, respectively; the minor-arc
exponents have still larger margins than required.
Lemma~\ref{kernel:coefficients} and $|V_N(y)|\le1$ now complete
the comparison on the major arcs as well. All thresholds used here
are uniform under $l\le N^\zeta$ and $D\le N^\beta$.
\end{proof}

\subsection{From Fourier comparison to forbidden differences}

\begin{proof}[Proof of Proposition~\ref{kernel:estimate}]
Extend $J_1,J_2$ by zero to $\Z$, and for either kernel define
\[
 \mathcal B_i=\frac1N\sum_{x,n\in\Z}k_i(n)J_1(x)J_2(x+n).
\]
Let $\widetilde J_i(\alpha)=\sum_xJ_i(x)e(-x\alpha)$ denote the
counting-measure Fourier transforms. Expanding the characters and
integrating imposes $n+x-y=0$, so
\[
 \mathcal B_i=\frac1N\int_0^1
 S_i(\alpha)\widetilde J_1(-\alpha)\widetilde J_2(\alpha)\,d\alpha.
\]
Cauchy--Schwarz, Plancherel and Lemma~\ref{kernel:comparison} give
\[
 |\mathcal B_1-\mathcal B_0|
 \le\frac{\|S_1-S_0\|_\infty}{N}
             \|J_1\|_{\ell^2(\Z)}\|J_2\|_{\ell^2(\Z)}
 \ll H^{-\gamma/2},
\]
because the two counting-measure norms are at most $\sqrt N$.

Every term in $\mathcal B_1$ uses a positive shift $F(m)$ with
$m\in\N_+$, and hence vanishes by avoidance. Thus
$\mathcal B_1=0$, regardless of the signed weights in~\eqref{eq:17}.
For a pair in the supports of $J_1,J_2$, ordering gives
$1\le y-x\le N-1$ and the residue conditions give
$y-x\equiv r\pmod D$. Both indicators in $k_0$ are therefore
automatic, and
\begin{align*}
 \mathcal B_0
 &=\frac D{N^2}\sum_{x,y\in[N]}J_1(x)J_2(y)
                    \sum_{\xi\in\Xi}\Gamma(\xi)e(-\xi(y-x))\\
 &=D\sum_{\xi\in\Xi}\Gamma(\xi)
                    \widehat{J_1}_{[N]}(-\xi)\widehat{J_2}_{[N]}(\xi).
\end{align*}
This proves~\eqref{eq:16} with its stated normalization and signs.
\end{proof}
